%------------------------------------------------------------------------------%
\section{Produtos notáveis}

Alguns produtos de expressões algébricas são encontrados tão frequentemente,
que são chamados de
\emph{produtos notáveis}\index{Produtos notaveis}.
Embora possamos calcular esses
produtos usando a propriedade distributiva, acabamos, com o uso, decorando as
fórmulas empregadas em sua obtenção. Segue abaixo alguns casos de produtos
notáveis. Em todos os casos considere que $a$ e $b$ são números reais,
variáveis ou expressões algébricas.

%------------------------------------------------------------------------------%
\emph{Quadrado da soma de dois termos}

A expressão $(a + b)^2$ representa a área de um quadrado de lado $(a+b)$.
Ao expandirmos, vemos que
\begin{align*}
  (a+b)^2
  & = (a+b) (a+b)           \\
  & = a^2 + ab + ba + b^2   \\
  & = a^2 + 2ab + b^2.
\end{align*}
Assim,
\begin{equation}
  \label{eq:produto_notavel_1}
  (a+b)^2 = a^2 + 2ab + b^2
\end{equation}
Regra Prática: O quadrado do primeiro termo, mais duas vezes o produto do
primeiro pelo segundo, mais o quadrado do segundo termo.

%------------------------------------------------------------------------------%
\emph{Quadrado da diferença de dois termos}

Funciona de forma idêntica à soma,
mas com atenção ao sinal negativo no termo central
\begin{align*}
  (a-b)^2
  & = (a-b) (a-b)         \\
  & = a^2 - ab - ba + b^2 \\
  & = a^2 - 2ab + b^2.
\end{align*}
Assim,
\begin{equation}
  \label{eq:produto_notavel_2}
  (a-b)^2 = a^2 - 2ab + b^2
\end{equation}
Regra Prática: O quadrado do primeiro termo, menos duas vezes o
produto do primeiro pelo segundo, mais o quadrado do segundo termo.

%------------------------------------------------------------------------------%
\emph{Produto da soma pela diferença}

Este é um dos mais úteis para simplificações e fatoração,
pois os termos centrais se anulam
\begin{align*}
  (a+b)(a-b)
  & = a^2 - ab + ba - b^2 \\
  & = a^2 - b^2.
\end{align*}
Assim,
\begin{equation}
  \label{eq:produto_notavel_3}
  (a+b)(a-b) = a^2 - b^2
\end{equation}
Regra Prática: Resulta sempre na diferença entre os quadrados dos dois termos.

%------------------------------------------------------------------------------%
\emph{Cubo da soma de dois termos}

Para desenvolvermos o cubo da soma $(a+b)^3$, seguimos os seguintes passos.
Escrevemos como o produto do quadrado pela base
\[
  (a+b)^3 = (a+b)^2 (a+b).
\]
Substituímos o resultado já conhecido do quadrado da soma
\[
  (a+b)^3 = (a^2 + 2ab + b^2) (a+b).
\]
Aplicamos a distributiva termo a termo
\[
  (a+b)^3 = a^3 + a^2b + 2a^2b + 2ab^2 + ab^2 + b^3
\]
Somamos os termos semelhantes para chegar na fórmula final
\begin{equation}
  \label{eq:produto_notavel_4}
  (a+b)^3 = a^3 + 3a^2b + 3ab^2 + b^3
\end{equation}

%------------------------------------------------------------------------------%
\emph{Cubo da diferença de dois termos}

\begin{equation}
  \label{eq:produto_notavel_5}
  (a-b)^3 = a^3 - 3a^2b + 3ab^2 - b^3
\end{equation}

%------------------------------------------------------------------------------%
\begin{example}
  Calcule o produto notável \((2x+5)^2\)

  \solution
  Nesse caso temos o quadrado da soma de dois termos
  \eqref{eq:produto_notavel_1}, então
  \begin{align*}
    (2x+5)^2
    & = (2x)^2 + 2\cdot2 x 5 + 5^2   \\
    & = 4x^2 + 20x + 25
  \end{align*}
\end{example}
%------------------------------------------------------------------------------%

%------------------------------------------------------------------------------%
\begin{example}
  Calcule o produto notável \((x-2)(x+2)\)

  \solution
  Aqui temos o produto da soma pela diferença
  \eqref{eq:produto_notavel_3}, logo
  \begin{align*}
    (x-2)(x+2)
    & = x^2 - 2^2   \\
    & = x^2 - 4
  \end{align*}
\end{example}
%------------------------------------------------------------------------------%

%------------------------------------------------------------------------------%
\begin{example}
  Calcule o produto notável \((y+4)^3\)

  \solution
  Usando a fórmula do cubo da soma de dois termos
  \eqref{eq:produto_notavel_4},
  \begin{align*}
    (y+4)^3
    & = y^3 + 3 y^2 4 + 3 y 4^2 + 4^3 \\
    & = y^3 + 12y^2 + 48y + 64
  \end{align*}
\end{example}
%------------------------------------------------------------------------------%

%------------------------------------------------------------------------------%
\begin{example}
  Calcule o produto notável \(\left(5-\frac{4}{y}\right)^2\)

  \solution
  Aqui temos o quadrado da diferença de dois termos
  \eqref{eq:produto_notavel_2}, portanto
  \begin{align*}
    \left( 5 - \frac{4}{y} \right)^2
    & = 5^2 - 2 \cdot 5 \frac{4}{y} + \left(\frac{4}{y}\right)^2   \\[2mm]
    & = 25 - \frac{40}{y} + \frac{16}{y^2}.
  \end{align*}
\end{example}
%------------------------------------------------------------------------------%

%------------------------------------------------------------------------------%
\begin{example}
  Calcule o produto notável \(\left(\sqrt[3]{x}-1\right)^3\)

  \solution
  Nesse caso temos o cubo da diferença, então
  \begin{align*}
    \left( \sqrt[3]{x} - 1 \right)^3
    & = \left( \sqrt[3]{x} \right)^3
      - 3 \left( \sqrt[3]{x} \right)^2 1
      + 3 \sqrt[3]{x} \;1^2
      - 1^3                                        \\[2mm]
    & = x - 3\sqrt[3]{x^2}+  3 \sqrt[3]{x} - 1
  \end{align*}
\end{example}
%------------------------------------------------------------------------------%
